\section{A pre-registered survey design and analysis plan}\label{sec:design}

This section is the design and analysis plan that we will register, time-stamped and with the commit hash of the analysis code, before any respondent is interviewed \citep{nosek2018}. It fixes in advance the item bank structure, the groups for invariance testing, the list-experiment module, the sample and poststratification frame, the estimands, the threshold $\tau$, the decision rules and the model checks. Any departure will be reported as a deviation, with its reason. The estimators are those of \cref{sec:model,sec:identification,sec:sensitive,sec:smallarea}, implemented and tested on synthetic data in \cref{sec:simulation}. No survey has been fielded. The items described below are \emph{constructs and selection criteria}, not a validated instrument. Item wording will be developed through the pre-testing steps of \cref{sec:design-items} and registered only once it is final.

\subsection{Population, constructs and model}\label{sec:design-constructs}

The target population is the adult household population of one country. The plan is written for the United States because the poststratification frame of \cref{sec:design-sampling} is available there; \cref{sec:design-groups} describes the cross-national extension. The measurement model is the bifactor graded response model of \cref{sec:model}, with a general factor $G$ and specific factors $S$ (classic stereotypes) and $I$ (Israel-related attitudes) \citep{samejima1969,gibbons1992,reise2012}. Responses are on a five-point agreement scale, so $C_j=4$ for every attitude item. The primary model adds an acquiescence factor \citep{maydeuolivares2006}: a person-level tendency to agree, orthogonal to $G$, $S$ and $I$, with one loading common to all attitude items, entered with a positive sign for agree-keyed items and a negative sign for reverse-keyed items, which are scored in the direction of their construct. $\mathrm{ECV}_j$ is computed from the content loadings as in \cref{def:boundary}. The identification results of \cref{sec:identification} are stated without this factor; its effect on recovery was checked by simulation. Without such a factor, acquiescence loads on $G$ and moves the boundary estimands (\cref{sec:model-meaning}); in the simulation the factor removed that bias, whereas balancing the keying without it did not (\cref{sec:design-decisions}). The correlated-factors model, with factors $S$ and $I$ and no general factor, and a unidimensional model are fitted alongside it as pre-registered comparison models.

\subsection{Item bank}\label{sec:design-items}

\Cref{tab:design-items} lists the item blocks. Every block is described by construct; none of the content descriptions below is a questionnaire item.

\begin{table}[tb]
\centering
\caption{Item bank structure. Counts are targets for the main survey; the pilot fields about 1.5 times as many candidates per block. $\mathcal{S}$, $\mathcal{I}$ and $\mathcal{A}$ are the item sets of \cref{sec:model}.}\label{tab:design-items}
\footnotesize
\setlength{\tabcolsep}{4pt}
\renewcommand{\arraystretch}{1.15}
\begin{tabularx}{\textwidth}{@{}>{\raggedright\arraybackslash}p{2.6cm}>{\raggedright\arraybackslash}p{1.9cm}c>{\raggedright\arraybackslash}X@{}}
\toprule
Block & Model role & Items & Construct and content (generic) \\
\midrule
S: classic stereotypes & $\mathcal{S}$; loads on $G$ and $S$ & 10 & Statements attributing to Jews as a group excessive power or influence (economic, political, media), disloyalty, conspiratorial intent, or responsibility for historical misfortunes; statements that minimise or deny the Holocaust. Content categories follow those used in prior batteries (\cref{sec:bg-surveys}) without reusing their wording. \\
\addlinespace
ID: Israel, disputed categories & $\mathcal{I}$; loads on $G$ and $I$ & 8 & Statements in the categories of Israel-related speech on which the definitions of \cref{sec:bg-definitions} disagree: positions on specific policies or actions of Israeli governments and of the Israeli state, on Zionism, on boycotts, and on comparisons with other historical cases. Selected to span those categories. \\
\addlinespace
IT: Israel, classic tropes & $\mathcal{I}$; loads on $G$ and $I$ & 6 & Israel-related statements that invoke classic antisemitic tropes, for example claims of secret control by Israel or its supporters, or collective responsibility of Jews for Israel's conduct. Selected from categories that more than one definition treats as antisemitic. \\
\addlinespace
A: anchors & $\mathcal{A}\subset\mathcal{S}\cup\mathcal{I}$ & 6 & Items expected, on substantive grounds and pilot evidence, to function equivalently across the groups of \cref{tab:design-groups}: four from S and two from the Israel-related blocks, so that every factor has anchors. Counted within their own blocks. \\
\addlinespace
C: comparison states & auxiliary & 16 & For each ID item, parallel statements about a comparable policy or case of each of two pre-registered comparison states. Not in the measurement model. \\
\addlinespace
R: reverse-keyed & $\mathcal{S}\cup\mathcal{I}$ & 4 & Favourably worded items, two in S and two in ID, which separate the acquiescence factor of the primary model from content (\cref{sec:design-constructs}). Counted within their own blocks. \\
\addlinespace
Q: data quality & auxiliary & 3 & Instructed-response attention checks, placed in separate blocks. \\
\addlinespace
K: knowledge and covariates & auxiliary & -- & Factual knowledge of the Israeli--Palestinian conflict and of other conflicts; ideology; religiosity and affiliation; birth year; region; demographics for poststratification. \\
\bottomrule
\end{tabularx}
\end{table}

\paragraph{Selection criteria}
Candidate items enter the pilot only if they meet five criteria.
\begin{enumerate}[(i)]
  \item \emph{One construct per item.} Each attitude item is written to load on $G$ and on exactly one specific factor. Double-barrelled statements (for example, a statement about a policy combined with a trope) are excluded, because they would make the item-level estimand of \cref{def:boundary} uninterpretable.
  \item \emph{Coverage of the disputed categories.} The ID block includes at least one item for each category of Israel-related speech on which \cref{sec:bg-definitions} records disagreement. The IT block draws on categories that more than one definition treats as antisemitic. Neither block is chosen to make the boundary set larger or smaller.
  \item \emph{Definitional tags, recorded but not used in estimation.} Before fielding, each Israel-related item is tagged with how the text of each definition treats its content category, or ``unclear''. Tags are assigned independently by advisory-group members who hold different views of the definitions (\cref{sec:ethics}), and disagreements are recorded rather than resolved. Tags are metadata. They let readers relate the estimated boundary set to any definition, but no estimand depends on them.
  \item \emph{Pre-testing.} Cognitive interviews with respondents who differ in ideology, age and religious background check that each item is understood as intended. Items that are understood differently by different groups are revised or dropped before the pilot; this is the cheapest stage at which to remove DIF.
  \item \emph{Identification.} \Cref{thm:ident} needs at least four items in one domain whose general-factor shares are not all equal save one, and the joint tetrad test of rule D2(ii) in \cref{sec:design-decisions} needs no more (its intersection--union diagnostic needs five). Both domains exceed these counts, with room to drop items at the pilot stage. The pilot also checks the loading conditions of \cref{thm:ident}, including heterogeneity of general-factor shares within a domain (\cref{asm:id-het}). A block whose specific-factor loadings are near zero in the pilot, or whose items have nearly equal general-factor shares, is redesigned before the main survey; the simulation shows that near-proportional loadings make the boundary estimands unreliable at realistic sample sizes (\cref{sec:sim-recovery}).
\end{enumerate}
\paragraph{Pilot}
A pilot sample, which is not used in any confirmatory analysis, screens candidates. Items are dropped for severe floor or ceiling effects, for high item non-response, or for DIF flagged in the pilot by both screening methods of \cref{sec:design-groups}. The pilot also supplies the variance inputs for \cref{thm:design}. The final item set, the anchor set $\mathcal{A}$ and all thresholds are frozen in the registration after the pilot and before the main fieldwork.

\paragraph{Administration}
The questionnaire is self-administered online, which tends to elicit more reporting of sensitive information than interviewer administration \citep{kreuter2008}. The list-experiment module (\cref{sec:design-list}) comes first, before any direct question about Jews or Israel. The order of the S and Israel-related blocks is randomised, and items are randomised within blocks. Block order is recorded and treated as a grouping variable for the invariance checks, so that priming between the two domains can be detected.

\subsection{Groups for invariance testing}\label{sec:design-groups}

\Cref{tab:design-groups} lists the pre-registered groupings. Each is tested separately; intersections are not tested, because their cells are too small.

\begin{table}[tbp]
\centering
\caption{Pre-registered groupings for invariance and DIF testing. Groups with fewer than $n_{\min}=300$ respondents are merged with an adjacent group or omitted, as fixed in the registration.}\label{tab:design-groups}
\footnotesize
\renewcommand{\arraystretch}{1.15}
\begin{tabularx}{\textwidth}{@{}l>{\raggedright\arraybackslash}X>{\raggedright\arraybackslash}X@{}}
\toprule
Grouping & Categories & Rationale \\
\midrule
Ideology & Self-placement, collapsed to left / centre / right & Political identity may shape responses to Israel-related items, and to some stereotype items, independently of $G$ (\cref{sec:bg-invariance}) \\
Cohort & Ages 18--29, 30--44, 45--64, 65+ & Generational differences in the salience and framing of both domains \\
Religiosity & Attendance at religious services: never; less than monthly; monthly or more & Religious commitments may affect both domains \\
Region & Four census regions & Required for the small-area model (\cref{sec:design-sampling}) \\
Block order & S first / Israel-related first & Detects priming between domains \\
Country (extension) & Countries in a cross-national replication & Translation and context; anchoring vignettes considered \citep{king2004} \\
\bottomrule
\end{tabularx}
\end{table}

For each grouping the procedure is fixed.
\begin{enumerate}[1.]
  \item Fit the configural multiple-group bifactor model, then constrain loadings (metric) and then thresholds (scalar) \citep{meredith1993,millsap2011}. Changes in fit are reported with the caveats of \citet{chen2007}; they do not by themselves decide invariance.
  \item Screen every item for DIF with two methods: Wald tests of equal slopes and thresholds in the multiple-group fit of the primary estimator, with the anchor set held invariant, the limited-information counterpart of the likelihood-ratio tests of \citet{thissen1993} and the test used in the simulation (\cref{sec:sim-dif}); and the Mantel--Haenszel procedure with the total score on the remaining items as the matching variable \citep{holland1988}. Multiplicity is controlled by the Benjamini--Hochberg procedure at a false discovery rate of 0.05 within each grouping \citep{benjamini1995}. An item is flagged if its Wald test is significant after this adjustment. The Mantel--Haenszel procedure, which is sensitive mainly to uniform DIF, is a cross-check, and disagreements between the two methods are reported. The size of each flagged item's DIF is a secondary screen. We report the difference in expected item score between the group-specific and the reference item parameters, averaged in absolute value over the group's latent distribution, and describe flagged items with a difference below 0.25 scale points as having small DIF; they are freed all the same. An earlier draft flagged an item only if it was also at least 0.25 scale points apart at the group's mean of $G$. In the simulation that requirement almost never flagged an item with non-uniform DIF, whose expected scores differ little near the group mean, and the procedure then reproduced the bias it was meant to remove (\cref{sec:sim-dif}). The revised screen was re-simulated end to end on the same data sets (\cref{sec:simulation}). With the correct anchors it freed the item with non-uniform DIF in $41\%$ of replications at the smallest DIF and in every replication at larger DIF; that item's focal ECV was then biased by $+0.042$ at the smallest DIF and by at most $0.002$ in absolute value at larger DIF, and the bias of the difference in means of $G$ was at most $0.006$ in absolute value.
  \item Fit the partial-invariance model with flagged items freed and the anchors constrained \citep{byrne1989}. Group comparisons of latent means and of $\mathrm{ECV}_j$ are reported from this model. Their identification rests on \cref{prop:anchor}. Comparisons of $G$ and of the group-specific ECVs are also reported, for every grouping, from a model anchored in the four stereotype anchors alone (\cref{rem:id-anchor-S}). That model need not assume that any Israel-related item functions identically across groups; it assumes that the stereotype anchors do, which is testable only in part (\cref{prop:anchor}(b)). It is a sensitivity analysis, not the primary comparison. Resting on four anchors instead of six, it is more exposed to DIF on them: in the simulation, DIF of $0.1$ on two stereotype anchors biased its difference in means of $G$ by $+0.073$ and $+0.080$, against $+0.025$ and $+0.031$ with the default anchors, and at $0.3$ it broke down (\cref{sec:sim-dif}). Every anchored comparison of latent means is therefore reported with the bound \eqref{eq:id-anchor-bound}, evaluated at the estimated information and maximised over all uniform anchor DIF of at most $\delta=0.2$ standard deviations of each anchor's latent response, the smallest Israel-item DIF in the simulation (\cref{sec:sim-dif}); the bound at $0.4$ is reported alongside. The bound is wide. In the simulated design, at $\delta=0.2$ it allows a first-order bias of the difference in means of $G$ of up to $0.69$ with the default anchors and $0.89$ with stereotype-only anchors, against a true difference of $0.3$; if all anchors' DIF has the same sign, the largest bias is about $0.2$ for both. As a further check on the anchors, each anchor is in turn freed and the comparisons recomputed; the spread of these leave-one-anchor-out estimates is reported.
  \item As a sensitivity analysis, refit with the alignment method \citep{asparouhov2014}, which does not require a pre-selected anchor set.
\end{enumerate}
Items flagged for non-uniform DIF have group-specific loadings, so their $\mathrm{ECV}_j$ is reported by group (\cref{sec:design-estimands}). \Cref{thm:dif-bias} gives the bias that the analysis would incur if it ignored flagged DIF; we report that bias, evaluated at the estimates, as a diagnostic.

\subsection{The list-experiment module}\label{sec:design-list}

\paragraph{Sensitive items}
The module covers four sensitive items drawn from the final item bank: two from S and two from the Israel-related blocks (one ID, one IT). For each sensitive item $j$, every respondent is randomised independently to one of the three arms of \cref{sec:sensitive}: \emph{direct} ($\mathsf{D}$), in which item $j$ appears in the direct battery; \emph{list control} ($\mathsf{C}$); or \emph{list treatment} ($\mathsf{T}$). Respondents in the two list arms are not asked item $j$ directly in the main questionnaire. For them, item $j$ is missing by design, and since assignment is random the missingness is ignorable in the measurement model. The prevalence $\pi_j$ targeted by the list is agreement with item $j$, coded as the top two categories of the five-point scale. The direct-arm counterpart is $\pi^D_j$, and $\mathrm{SDB}_j=\pi_j-\pi^D_j$ as in \cref{sec:sensitive}.

\paragraph{Control lists}
Each sensitive item has its own list of $J$ control items ($J$ as in \cref{sec:sensitive}). $J$ and the arm probabilities $(r_D,r_C,r_T)$ are chosen by \cref{thm:design}, with $J\in\{3,4,5\}$ and the inputs (control-item prevalences and their covariance) estimated in the pilot. Control items are low in sensitivity and unrelated to Jews or Israel. Following \citet{glynn2013}, each list includes negatively correlated pairs of control items. This reduces the variance of the count and makes it rare for a respondent to agree with all or none of the control items, which limits ceiling and floor effects. It can also make the count more informative about the sensitive answer (\cref{sec:sens-design}), so every list must meet two registered disclosure bounds, evaluated at the pilot estimates of the control-count distribution and of the sensitive item's prevalence: the mutual information $I(Y_i;Z_i\mid\mathsf{T})$ between the treatment count and the sensitive answer is at most $0.15$ bits, and at most $15\%$ of treated respondents have posterior $\Pr(Z_i=1\mid Y_i,\mathsf{T})\ge0.5$. For comparison, four independent control items of prevalence $1/2$ with $\pi=0.2$ give $0.11$ bits and $11\%$, and two pairs with within-pair correlation $-0.6$ give $0.24$ bits and $30\%$. Within the lists that meet both bounds, $J$ and the correlation are chosen by \cref{thm:design}(d); a list that fails either bound is replaced before the main survey.

\paragraph{Checks}
For each list we report the test of \citet{blair2012} for design effects. The decision rule is in \cref{sec:design-decisions}. We also estimate the floor-and-ceiling model of \citet{blair2012}, and report the multivariate estimator of \citet{imai2011} with the covariates of \cref{tab:design-groups}. In addition, a random 20\% of list-arm respondents are asked item $j$ directly at the very end of the questionnaire. This supports the combined estimator and placebo tests of \citet{aronow2015} as a secondary analysis. It does not enter the primary estimator, because a direct answer given after the list may be affected by it. For these respondents the file records the direct answer as well as the count, so the list gives them no additional protection; \cref{sec:ethics-data} sets out how these records are handled.

\paragraph{The cost of the module}
\Cref{prop:list-var} gives the variance of the difference-in-means estimator. As illustrative arithmetic, not a result, suppose that both count variances equal $1$ ($\sigma_{\mathsf{C}}^2=\sigma_{\mathsf{T}}^2=1$), that $\pi^D_j=0.2$, and that the three arms are of equal size $m$; these are the inputs used for \cref{thm:design}(a) in \cref{app:pd-design}. Equal count variances assume that the control items covary slightly negatively with the sensitive item. With four independent control items of prevalence $1/2$ and $\pi_j=0.2$, $\sigma_{\mathsf{T}}^2=1.16$ instead (\cref{sec:sens-var}). Under the first assumption the variance of $\widehat{\mathrm{SDB}}_j$ is about $(2+0.16)/m$, and detecting $\mathrm{SDB}_j=0.10$ with power 0.8 in a two-sided test at the 5\% level requires $m\approx 1{,}700$, or about 5,100 respondents in total for that item. With $\sigma_{\mathsf{T}}^2=1.16$ it requires $m\approx1{,}820$. Detecting $\mathrm{SDB}_j=0.05$ requires four times as many, about 20{,}000 respondents for each item, against a planned 6{,}000. A meta-analysis of list experiments found sensitivity biases typically below ten percentage points \citep{blair2020}, so the module is not powered for the effects that are most plausible. $\mathrm{SDB}_j$ is therefore a secondary estimand (E5 in \cref{sec:design-estimands}), with a registered precision target instead of a power target: at the planned sample size and the allocation of \cref{thm:design}, the 95\% interval for each $\mathrm{SDB}_j$ should have a half-width of at most $0.07$ (about $0.06$ under the inputs above). Two features of the design reduce the cost. Each respondent serves in some arm for every sensitive item, so the four experiments share one sample. And the joint model of \cref{thm:joint-sdb} borrows strength across items and from the measurement model. Equal allocation is not optimal; \cref{thm:design} gives the allocation actually used.

\subsection{Sampling, weighting and poststratification}\label{sec:design-sampling}

\paragraph{Primary sample}
The primary sample is drawn from a probability-based online panel recruited by address-based sampling, with a target of 6,000 completed interviews. The target will be revised in the registration from the power analysis of \cref{sec:design-power}. Design weights are raked to population margins for age, sex, education, race and ethnicity, and region. These weights are used for the national estimands.

\paragraph{Optional opt-in supplement}
For small-area estimation, a non-probability (opt-in) supplement may be added in under-sampled areas. Opt-in samples can be less accurate than probability samples, even after demographic weighting, and their accuracy varies widely across surveys and measures \citep{yeager2011,baker2013}. Model-based adjustment can nonetheless make them informative \citep{wang2015}. The supplement therefore enters only the MRP model. It carries an indicator for sample source, and its estimates are compared with those from the probability sample in areas where both are available. If the two disagree by more than the pre-registered tolerance (twice the combined standard error of an area mean of $G$), the supplement is dropped.

\paragraph{Poststratification}
The small-area estimand is the poststratified mean in \cref{eq:mrp} \citep{gelman1997,park2004}. Cells $k$ are formed by state $\times$ age (4) $\times$ sex (2) $\times$ education (4) $\times$ race and ethnicity (4). Population counts $N_k$ come from the most recent five-year American Community Survey tables available at registration. Ideology and religiosity are not in census tables. For estimates that condition on them, the cell table is extended with a joint distribution estimated from a large reference survey, and the uncertainty of that extension is propagated. State-level predictors are fixed at registration. Multilevel regression with poststratification performs well relative to disaggregation of national samples, especially when state samples are small \citep{lax2009}. Its limits are those of the model: the estimates depend on the predictors, and design-based and model-based weighting can diverge \citep{gelman2007}. We therefore report a design-based estimate alongside every MRP estimate for areas large enough to have one.

\paragraph{Campus module (optional)}
Campus-level estimates require a separate frame of students. If such a module is fielded, its poststrata are institution type $\times$ enrolment characteristics, with counts from institutional enrolment data. Estimates for named institutions will not be published (\cref{sec:ethics}).

\subsection{Pre-registered estimands}\label{sec:design-estimands}

\paragraph{Primary estimands}
\begin{enumerate}[(E1)]
  \item For every Israel-related item $j\in\mathcal{I}$, the item-level explained common variance $\mathrm{ECV}_j$ of \cref{def:boundary}, with a 95\% interval.
  \item The boundary set $\mathcal{B}_\tau=\{j\in\mathcal{I}:\mathrm{ECV}_j\ge\tau\}$ with $\tau=0.5$, the point at which the general factor accounts for at least half of an item's common variance. The value of $\tau$ is fixed here, before any data are seen. Each item is classified three ways from simultaneous 95\% intervals for all $\mathrm{ECV}_j$, $j\in\mathcal{I}$: Bonferroni-adjusted Wald intervals on the logit scale (\cref{prop:id-boundary-test}(b)). An item is \emph{in} $\mathcal{B}_\tau$ if its lower limit is at least $\tau$, \emph{out} if its upper limit is below $\tau$, and \emph{undetermined} otherwise. With simultaneous intervals, the items classified in and the items not classified out bracket $\mathcal{B}_\tau$ with probability at least 0.95 in large samples, which pointwise intervals do not guarantee; in the simulation they did so in every replication with six Israel-related items (\cref{sec:sim-recovery}). Undetermined is a reportable outcome, not a failure.
  \item The latent correlation between the stereotype factor and the Israel-attitude factor in the correlated-factors model, with a 95\% interval.
  \item Differences in the latent mean of $G$ between the groups of \cref{tab:design-groups}, under the partial-invariance model, with stereotype-only anchoring as a sensitivity analysis (\cref{sec:design-groups}).
\end{enumerate}

\paragraph{Secondary estimands}
\begin{enumerate}[(E1)]\setcounter{enumi}{4}
  \item For each sensitive item $j$, $\mathrm{SDB}_j$ from the joint model of \cref{thm:joint-sdb}, with the precision target of \cref{sec:design-list}.
  \item $\mathrm{ECV}_j$ and $\mathcal{B}_\tau$ by group, with the same simultaneous rule, for groupings in which loadings are not invariant.
  \item MRP estimates of the mean of $G$ by state and census region (\cref{eq:mrp}); boundary estimands are estimated by region only, as a multiple-group model.
  \item Within-person contrasts between each ID item and its comparison-state counterparts (block C), as a description of whether the same kind of statement is endorsed differently about Israel and about other states. A contrast is not by itself evidence of a double standard: respondents may know more about one case, so contrasts are also reported conditional on the knowledge measures of block K \citep{beattie2017}.
  \item The proportion agreeing with at least $q$ of the $M$ S items (\cref{sec:sa-index}), computed for comparability with published indices, and reported by group only together with the invariance diagnostics that \cref{prop:index} shows to be necessary.
\end{enumerate}
\paragraph{Sensitivity analyses}
$\mathcal{B}_\tau$ is also reported for $\tau\in\{0.3,0.7\}$; these do not change the primary classification. Other sensitivity analyses: exclusion of respondents who fail two or more attention checks (the primary analysis keeps them, because passing attention checks correlates with respondent characteristics; \citealp{berinsky2014}); a refit in which $\mathcal I$ is the ID block alone, which removes the trope-based items most exposed to a cross-loading on $S$; the primary model without the acquiescence factor; group comparisons anchored in the stereotype items alone, with the bound \eqref{eq:id-anchor-bound} (\cref{sec:design-groups}); the alignment model; the bifactor-$(S-1)$ model of \citet{eid2017}, in which the general factor is defined by the stereotype items alone (\cref{rem:id-s1}); percentile-bootstrap intervals with at least 1{,}000 resamples, as a check on the Wald intervals of E1 and E2; full-information and Bayesian fits (\cref{sec:design-decisions}); and Schmid--Leiman and exploratory bifactor solutions \citep{schmid1957,reise2012}.

\subsection{Model checks and decision rules}\label{sec:design-decisions}

\paragraph{Estimation and checks}
The primary estimator is the two-step limited-information estimator validated in the simulation (\cref{sec:sim-design}): thresholds and polychoric correlations \citep{olsson1979}, with an asymptotic covariance matrix from per-respondent influence functions, then weighted least squares with diagonal weights and sandwich standard errors, from a default start and at least 20 random starts. It is primary because it is the estimator whose recovery, coverage and identification diagnostic the simulation established; the diagnostic of rule D2(iii) and the power study of \cref{sec:design-power} are defined for it. Full-information marginal maximum likelihood, exploiting the dimension reduction of the bifactor structure \citep{gibbons1992}, and a Bayesian fit with weakly informative priors \citep{gelman2013} are reported as sensitivity analyses. Under the normal-ogive model all three target the same parameters (\cref{lem:id-reduction}), so a material disagreement between them is itself reported as a finding about the fit. Model adequacy is assessed with limited-information fit statistics, residual correlations between item pairs, and, for the Bayesian fit, posterior predictive checks of item-pair and sum-score distributions. The bifactor indices of \citet{rodriguez2016} are reported for all three models: explained common variance, $\omega_H$, percentage of uncontaminated correlations and construct replicability. Their thresholds are descriptive and do not enter the decision rules.

\paragraph{Cross-loadings and response style}
Two misspecifications would move the ECVs of Israel-related items, though not all in the same direction: a loading of trope-based items on the stereotype factor, which a bifactor model without that loading credits to $G$, and acquiescence (\cref{sec:model-meaning}). In the simulation (\cref{sec:sim-misspec}), a cross-loading of $0.3$ raised the mean ECV of the IT items by $0.089$ and lowered that of the ID items by $0.080$, and with the registered keying acquiescence moved the ECVs of single Israel-related items by between $+0.135$ and $-0.312$, upwards for agree-keyed and downwards for reverse-keyed items. Acquiescence is modelled in the primary model (\cref{sec:design-constructs}). In the simulation of both misspecifications at the registered item counts and $n=6{,}000$, the acquiescence factor kept the bias of every Israel-related ECV below $0.007$ under both the registered and a balanced keying, against biases of up to $0.43$ without it; balancing the keying without the factor did not remove the bias. With acquiescence slopes that differ between items, which the factor does not model, the largest bias was $0.043$, against $0.112$ without it. For cross-loadings, the mean residual polychoric correlation between the IT and the S items under the primary fit is tested one-sided at the 5\% level. If the test rejects, the ECVs and classifications of the Israel-related items are taken from a model in which every IT item also loads on $S$, with the cross-loadings of the whole IT block fitted jointly, $\mathrm{ECV}_j$ still computed from the loadings on $G$ and $I$, and the estimated cross-loadings reported. In the simulation, with the test computed from the bifactor fit, it rejected in $85$ to $100\%$ of replications when the IT items cross-loaded on $S$ with standardised loadings of $0.1$ to $0.3$ (in $5\%$ without cross-loadings), and this two-stage rule kept the bias of every Israel-related ECV below $0.021$, against biases of up to $0.141$ for the bifactor model alone. A check that fits a cross-loading for one IT item at a time, which an earlier draft proposed, is not used: at $c=0.3$ it cut the share of IT items on the wrong side of $\tau$ from $14.2\%$ to $1.5\%$ only by leaving $48.5\%$ of them undetermined, and its one-item fits stayed biased, because the other IT items' cross-loadings remain unmodelled (\cref{sec:sim-misspec}). The refit with $\mathcal I$ restricted to the ID block (\cref{sec:design-estimands}) is the corresponding check that does not rely on the test.

\paragraph{Decision rules}
The following rules are fixed in advance.
\begin{description}[style=unboxed,leftmargin=0pt]
  \item[D1, structure.] If the bifactor model fails the pre-registered fit criteria and the correlated-factors model passes them, we report E3 and no boundary set. If both fail, we report descriptive item statistics and the model comparison only.
  \item[D2, identification.] Item-level ECVs and the boundary set are reported only if the data support the conditions of \cref{thm:ident}. (i) Each domain has at least three items whose specific loadings have 95\% intervals that exclude zero (\cref{asm:id-spec}). (ii) In at least one domain $\mathcal D$ the items are heterogeneous in the sense of \cref{asm:id-het}. By \cref{lem:app-tetrad,lem:app-classes}, the assumption fails in $\mathcal D$ exactly when every polychoric tetrad $R_{k_1k_2}R_{k_3k_4}-R_{k_1k_3}R_{k_2k_4}$ among its items vanishes (under the bifactor model each equals $D_{k_1k_4}D_{k_2k_3}$). Its failure is therefore the null hypothesis of a single test, the simultaneous Wald test of a non-redundant set of $|\mathcal D|(|\mathcal D|-3)/2$ tetrads of $\mathcal D$ \citep{bollen1993tetrad}. It uses only the polychoric correlations and their asymptotic covariance, so it is regular whether or not the bifactor model is identified, and it replaces the Bonferroni rule over all quadruples of an earlier draft. Each domain is tested at the 2.5\% level, so that the rule, which needs heterogeneity in only one domain, has level at most 5\%. At $n=6{,}000$ and the registered item counts the simulation found sizes close to nominal and power of one in the heterogeneous designs (\cref{sec:simulation}); at a smaller sample size the critical values are calibrated by simulation, because the test over-rejects at $n=1{,}500$ in the larger domain. The joint test needs only four items in $\mathcal D$, which give two non-redundant tetrads. We also report the intersection--union version \citep{berger1982}, which accepts heterogeneity only if the same test rejects for $\mathcal D\setminus\{k\}$ for every $k\in\mathcal D$, and so needs five items. The null hypothesis for $\mathcal D\setminus\{k\}$, that all its tetrads vanish, says that all its items but at most one lie on one line, not that its loadings are proportional. The intersection--union version therefore tests a stronger property, heterogeneity that survives the removal of any one item, and is conservative; it is a robustness diagnostic and does not decide. Neither test detects near-proportional loadings, against which (iii) is the guard. (iii) The fit is numerically identified: the identification diagnostic of \cref{sec:sim-recovery}, the ratio of the smallest to the largest singular value of the weighted Jacobian of the fitted polychoric correlations of the primary fit with respect to its content loadings (on $G$, $S$ and $I$), exceeds a cut-off fixed in the registration from the simulation of \cref{sec:design-power}, and at least 20 random starts reach the same optimum with no loading at a bound. The slope of the acquiescence factor is left out of the diagnostic because it enters the correlations only through its square: without acquiescence its information vanishes and the diagnostic over all parameters is zero, whereas over the content loadings it was $0.082$ to $0.105$ in the simulation, close to the value for the model without the factor (\cref{sec:sim-misspec}). If (i) or (ii) fails, every $\mathrm{ECV}_j$ is reported as not identified, no boundary set is reported, and E3 is reported instead; if only two Israel-related items carry specific variance, the identified sets of \cref{thm:ident}(c) are reported. If only (iii) fails, the ECVs are reported as weakly identified and every item is classified as undetermined. The interval-based classification of E2 is used only when (iii) holds, because the guarantees of \cref{prop:id-boundary-test} are pointwise and unreliable near the non-identified set (\cref{rem:id-test-uniform}). Agreement across starts is necessary, not sufficient: in the simulation it exposed the two-item case but neither the proportional nor the near-proportional case (\cref{sec:sim-recovery}).
  \item[D3, invariance.] Group comparisons (E4, E6) are reported only for groupings in which the anchors that pass the DIF screen satisfy the rank condition of \cref{prop:anchor}(a). There must be at least two such anchors in each domain, at least one with a non-zero specific loading in each domain, and, within at least one domain, two anchors whose loading vectors are not proportional, judged by a 95\% interval for their determinant $D_{kl}$ \eqref{eq:id-det} that excludes zero. Two proportional anchors per domain do not suffice. If no Israel-related anchor passes but the stereotype anchors contain a non-proportional pair, comparisons of $G$ and of the group-specific ECVs are reported, labelled as resting on the stereotype anchors alone, and comparisons of the Israel-attitude mean are not (\cref{rem:id-anchor-S}). Otherwise only within-group estimates and the alignment results are reported, labelled exploratory (\cref{prop:anchor}). Whatever the outcome, the stereotype-anchored comparisons and the bound \eqref{eq:id-anchor-bound} of step~3 in \cref{sec:design-groups} are reported alongside.
  \item[D4, design effects.] If the design-effect test rejects at the 5\% level for a list, that list's estimate is not used to correct the direct estimate; E5 for that item is reported as not estimable. If floor or ceiling effects are detected, the floor-and-ceiling model of \citet{blair2012} replaces the difference in means.
  \item[D5, completeness.] Every primary estimand is reported for every item and group, whatever its direction or magnitude. Items are dropped after fielding only for the pre-registered quality reasons (item non-response above 10\% or failed comprehension in the pilot), and dropped items are listed.
\end{description}

\subsection{Power and precision}\label{sec:design-power}

The sample size is set by Monte Carlo simulation, a practical approach to power analysis for latent-variable models of this size \citep{muthen2002}. We use the simulation code of \cref{sec:simulation}, which implements the primary estimator, with item parameters set from the pilot. The registration will report three precision targets: the half-width of the 95\% interval for each $\mathrm{ECV}_j$ at the national level; the power to detect DIF with an expected-score difference of 0.25 scale points in the smallest group of \cref{tab:design-groups}; and the half-width of the 95\% interval for each $\mathrm{SDB}_j$ under the allocation of \cref{thm:design} (target $0.07$; \cref{sec:design-list}). The same simulation fixes the cut-off of the identification diagnostic in rule D2(iii) and reports the power of the joint tetrad test of rule D2(ii), and of its intersection--union version, at the pilot parameters. The recovery results of \cref{sec:simulation} indicate which of these binds. If the budget cannot meet all three, the list module is reduced (fewer sensitive items) before the main sample is, because E1--E4 do not depend on it.

\subsection{Replication package and transparency}\label{sec:design-replication}

The following will be released. The registration document, with the commit hash of the analysis code. The code for every estimator, the simulation study and the figures. The final item bank, including wording, block, definitional tags and the record of tag disagreements, released after fieldwork closes. Aggregate results for every estimand, and the poststratification table. A log of deviations from the plan. De-identified microdata will be deposited with a data archive under restricted access, because the data combine religion, political opinion and geography (\cref{sec:ethics}). Synthetic microdata generated from the fitted model will be released openly, so that the code can be run without access to the restricted file.

\begin{table}[tbp]
\centering
\caption{Timeline, in months from the start of the project.}\label{tab:design-timeline}
\small
\begin{tabular}{@{}lll@{}}
\toprule
Months & Phase & Output \\
\midrule
1--2 & Advisory group convened; candidate items drafted; tagging & Candidate bank; tag record \\
3--4 & Cognitive interviews; ethics review & Revised bank; approval \\
5 & Pilot fieldwork and screening & Pilot report (not confirmatory) \\
6 & Item bank, anchors, $J$, allocation and thresholds frozen & Registration, with code hash \\
7--8 & Main fieldwork & Raw data, held by panel provider \\
9 & Data-quality checks; data freeze & Frozen analysis file \\
10--11 & Confirmatory analysis as registered & Results for E1--E9; deviations log \\
12 & Release & Paper, replication package, data deposit \\
\bottomrule
\end{tabular}
\end{table}

\Cref{tab:design-timeline} gives the timeline. Fieldwork dates will be recorded and reported with every estimate. Attitudes about Israel may move with events, so the estimates describe the population at the time of fieldwork and are not assumed to be stable.
